\section*{序論}
10月29日朝までの生成AI関連Xでは、Google DeepMindによるGemini 4 Carbonの発表が大きな話題となった。サイバーセキュリティと長時間にわたるタスク処理の性能を訴求する一方、当初の提供対象を限定する方針も示され、フロンティアモデルの能力と運用上の制御が並んで注目された。前日に発表された世界モデルGenie 4についても、デモや展開先を巡る議論が続いた。

公開モデルの領域では、DiffuSpaceの拡散言語モデルCloverとLiquid AIの小型モデルLFM3が取り上げられた。生成方式や推論効率、端末上での実行可能性が異なる方向から模索されている。また、Articul8のBiviaやVangridに関する投稿では、AIを実際の業務判断や物理世界のデータ取得につなぐ動きが語られた。

本号はGrokが収集した8件を原文の順序で整理する。AmazonのNova再編やDeepSeek V5に関する投稿には未確認情報が含まれる。さらに一部の投稿時刻は観測窓を数分越えており、収集内容と厳密な観測範囲の境界を明示する。独立した一次資料の再調査は行っていない。
